\icode{E-MD-SI-L3-261114-002}
量纲分析是我们分析、估算物理问题的有力工具和重要思路。1950年，英国力学家泰勒通过原子弹爆炸后火球随时间扩散的照片，并结合量纲分析，估算出了美国第一颗原子弹爆炸释放的能量。假设原子弹爆炸时形成的冲击波是球面波，爆炸中心是该球面波的球心，爆炸火球半径 \(R\) 仅仅依赖于爆炸后的时间 \(t\)，爆炸瞬间释放的能量 \(E\)、空气密度 \(\rho\)，以及无量纲常数 \(C\)，因此可以写成 \(R = C t^x E^y \rho^z\)。为了估算量级，\(C\) 可以近似等于1，空气密度 \(\rho = 1.2 \, \text{kg/m}^3\)。当原子弹爆炸25毫秒时刻的冲击波半径\(R =\) \,\qty{132}{m}。请你利用所学知识，估算该次原子弹爆炸释放的能量量级最接近以下哪个选项
\begin{choices}
	\item \qty{e8}{J}
	\item \qty{e14}{J}
	\item \qty{e20}{J}
	\item \qty{e25}{J}
\end{choices}
\begin{answer}
    B
\end{answer}